% Figaro Paper C: The Figaro Kernel — A Verified Solidity Implementation
% Implementation + verification snapshot. The mechanism is presented in
% the companion paper figaro3a.tex; organizational implications appear
% in figaro3b.tex.
% Snapshot: 2026-04-23

\documentclass[11pt,a4paper]{article}

% ── Packages ────────────────────────────────────────────────────────
\usepackage[utf8]{inputenc}
\usepackage[T1]{fontenc}
\usepackage{amsmath,amssymb,amsthm}
\usepackage{mathtools}
\usepackage{booktabs}
\usepackage{graphicx}
\usepackage{hyperref}
\usepackage{cleveref}
\usepackage{enumitem}
\usepackage{xcolor}
\usepackage{listings}
\usepackage[margin=1in]{geometry}
\usepackage{natbib}
\bibliographystyle{plainnat}

% ── Theorem environments ────────────────────────────────────────────
\newtheorem{theorem}{Theorem}[section]
\newtheorem{proposition}[theorem]{Proposition}
\newtheorem{corollary}[theorem]{Corollary}
\theoremstyle{definition}
\newtheorem{definition}[theorem]{Definition}
\theoremstyle{remark}
\newtheorem{remark}[theorem]{Remark}

% ── Listings config ─────────────────────────────────────────────────
\lstset{
  basicstyle=\ttfamily\small,
  keywordstyle=\bfseries\color{blue!70!black},
  commentstyle=\itshape\color{gray},
  frame=single,
  breaklines=true,
  captionpos=b,
}

% ── Macros ──────────────────────────────────────────────────────────
\newcommand{\figaro}{\textsc{Figaro}}
\newcommand{\pay}{P}
\newcommand{\cumval}{G}
\newcommand{\bbond}{C_b}
\newcommand{\sbond}{C_s}

% ────────────────────────────────────────────────────────────────────

\title{%
  The \figaro{} Kernel:\\
  A Verified Solidity Implementation of Asymmetric Bonding%
}

\author{
  Alessandro Daliana\thanks{Corresponding author. \texttt{adaliana@figaro.org}}
}

\date{April 2026 (snapshot: 2026-04-23)}

\begin{document}

\maketitle

% ════════════════════════════════════════════════════════════════════
\begin{abstract}
We describe a reference Solidity implementation of the
asymmetric-bonding settlement mechanism developed
in~\citep{daliana2026mechanism}, together with the formal-verification
stack used to audit it. The kernel (\texttt{FigaroCore}) is
\textbf{ownerless, fee-less, and admin-less}: two external functions,
three storage mappings, no upgrade path, no escape hatch from the
bonded state. Verification is layered: TLA\textsuperscript{+} model
checking exhaustively explores 6{,}087{,}113 distinct states across
seven kernel invariants; Echidna fuzzes seven property invariants
under a 50{,}000-call campaign budget; Halmos symbolically proves
seven kernel properties via the z3 SMT solver; Certora cloud proves
twenty-four declared CVL rules (twenty-five sub-rules in the cloud
report) across the kernel, the attestation surface, and a
token-operations conservation surface that gates every ERC-20 call
site in the source tree.

The implementation also realizes a \emph{coordinator pattern}: an
extension discipline under which external mechanisms (Dutch auction,
schema-typed attestation coordinator, schema and operator registries)
compose with the kernel without weakening the bonding equilibrium
of~\citep{daliana2026mechanism}. We give the four sufficient
conditions and three concrete instances. Finally, we describe the
three-layer enforcement architecture---economic (bonding),
coordinational (atomic resolution), evidentiary (immutable
log)---and the threat model under which the layers compose.

This paper is a snapshot at commit-of-record 2026-04-23. Verification
counts will grow as the test surface expands; the methodology is the
stable artifact.
\end{abstract}

\medskip
\noindent\textbf{Keywords:}
smart contracts, formal verification, TLA\textsuperscript{+}, Halmos,
Certora, Echidna, EIP-712, settlement layer, coordinator pattern

% ════════════════════════════════════════════════════════════════════
\section{Introduction}
\label{sec:intro}

The asymmetric-bonding mechanism developed
in~\citep{daliana2026mechanism} establishes that, given a settlement
layer with the operations \texttt{commit} and \texttt{resolveProcess},
cooperation is the unique outcome surviving iterated elimination of
weakly dominated strategies in an $N$-party process tree. The
mechanism is implementation-agnostic. This paper presents one
implementation of the mechanism---a Solidity smart contract
\texttt{FigaroCore} on the EVM---and the verification methodology
applied to it.

The core property the verification is asked to deliver is precisely
the gap between mechanism and code: the mechanism's theorems hold
\emph{by assumption} that the settlement layer enforces buyer
dominance, monotonic accumulation, and atomic resolution. The code
must \emph{actually} enforce them, in every reachable path, against
all reasonable adversaries. Bridging that gap is the work of the
verification stack.

\paragraph{What verification does and does not transfer from the mechanism paper.}
A distinction worth naming up front: the companion mechanism paper
proves \emph{equilibrium} properties of a bonding game --- that
cooperation weakly dominates defection at every position in an
$N$-party process tree, under specified rationality assumptions.
The present paper verifies \emph{structural preconditions} that
those equilibrium properties assume --- buyer dominance, monotonic
accumulation, atomic resolution, conservation of value. These are
different kinds of claim. The equilibrium itself is a property of
rational play over a payoff structure; it is not a property of
code and it is not verified by any of the four tools we apply.
What we verify is that the code implements the payoff structure
faithfully in every reachable path. When the companion
cryptoeconomics paper~\citep{daliana2026fig} asserts that a
token-emission policy ``preserves the kernel's bonding
equilibrium,'' it is invoking the mechanism paper's equilibrium
claim (which holds by rational play) and not the verification
claim made here (which holds by code inspection). Keeping the two
claims distinct matters for any reader who wants to know exactly
which property is backed by which evidence.

\paragraph{Brief recap of the mechanism.}
For a transaction with payment $\pay$ and cumulative upstream value
$\cumval \geq \pay$, the buyer locks $2\pay$ and the seller locks
$2\cumval$. Only the buyer can trigger resolution; on resolution the
buyer recovers $\pay$ and the seller recovers $2\cumval + \pay$. All
orders within a process resolve atomically. The full game-theoretic
treatment, including the two-party Nash result, the $N$-party
generalization, the escape-hatch impossibility, and the reduction
over joint-liability microfinance, is in~\citep{daliana2026mechanism}.

\paragraph{Paper organization.}
\Cref{sec:architecture} describes the contract architecture and the
EIP-712 commitment protocol. \Cref{sec:non-decisions} catalogues the
features deliberately absent from the kernel and ties each absence
to a mechanism property it would have broken.
\Cref{sec:verification-method} describes the four-layer verification
methodology. \Cref{sec:verification-results} presents the current
results. \Cref{sec:security} develops the threat model and
three-layer enforcement architecture.
\Cref{sec:composition} formulates the coordinator pattern and
exhibits three concrete instances. \Cref{sec:conclusion} concludes
with notes on agent-native operation and execution-environment
portability.

% ════════════════════════════════════════════════════════════════════
\section{Smart Contract Architecture}
\label{sec:architecture}

The kernel is a single Solidity 0.8.26 contract,
\texttt{src/FigaroCore.sol}, inheriting OpenZeppelin's \texttt{EIP712}
and \texttt{ReentrancyGuard}. It exposes two external functions:
\texttt{commit} and \texttt{resolveProcess}. There is no constructor
parameter, no owner, no admin, no upgrade path.

\subsection{Storage Layout}

Three mappings comprise the kernel state:
\begin{itemize}
  \item \texttt{processes}: \texttt{bytes32} $\to$
    \texttt{ProcessState}, where \texttt{ProcessState} is a struct of
    \texttt{(address rootBuyer, IERC20 currency, uint256
    cumulativeValue, uint256 activeOrderCount)}.
  \item \texttt{orderStatus}: \texttt{bytes32} $\to$ \texttt{uint8}
    encoding $\{0:\text{unknown}, 1:\text{committed},
    2:\text{resolved}\}$. The status function is monotonically
    non-decreasing; this is one of the verified invariants.
  \item \texttt{orderProcessId}: \texttt{bytes32} $\to$
    \texttt{bytes32} recording each order's process membership.
\end{itemize}

The kernel deliberately does not store a per-order struct. Order
terms (payment, cumulative-value snapshot, parties, denomination) are
reconstructed off-chain from the signed commitment and the emitted
\texttt{OrderCommitted} event. Bond amounts ($2\pay$ for the buyer,
$2\cumval$ for the seller) are deterministic functions of those
values rather than separate storage variables. Storage is therefore
$O(\text{processes}) + O(\text{orders})$ uint8 + $O(\text{orders})$
bytes32, with no order-payload duplication.

\subsection{Commitment Protocol}

Both parties sign an EIP-712 typed commitment~\citep{eip712}
off-chain. A single on-chain \texttt{commit} call verifies both
signatures (via \texttt{ECDSA.recover} against the EIP-712 digest)
and pulls both bonds atomically. Root commitments create a new
process; sub-order commitments extend an existing one by carrying
the inherited \texttt{processId} and the next expected cumulative
value. This eliminates the accept-reject pattern of traditional
escrow protocols (which front-run on acceptance) and makes the act
of committing simultaneous from the chain's point of view.

\subsection{Identifier Derivation}

Process and order identifiers are content-addressed:
\begin{align}
  \text{processId} &= H_{\text{EIP-712}}(\text{root commitment}), \\
  \text{orderHash} &= \mathrm{keccak256}(\text{processId} \;\|\; \text{structHash}).
\end{align}
The EIP-712 domain separator binds the root identifier to chain id
and verifying contract, preventing replay across deployments. Sub-orders
inherit \texttt{processId} from the signed commitment. No
auto-incrementing counter exists; identifier collisions require
SHA-3 collisions.

\subsection{Fee-on-Transfer Protection}

The internal helper \texttt{\_pullExact(IERC20 token, address from,
uint256 amount)} reads the contract's pre-transfer balance, calls
\texttt{safeTransferFrom}, and reverts with
\texttt{FeeOnTransferDetected} if
$\texttt{balanceAfter} - \texttt{balanceBefore} \neq
\texttt{amount}$. This ensures the conservation law cannot be broken
by tokens that apply a transfer tax---a class of issues that has
silently corrupted accounting in several DeFi systems.

\subsection{Events}

Five events are emitted, supporting full off-chain state
reconstruction without a subgraph dependency:

\begin{itemize}
  \item \texttt{OrderCommitted} carries the canonical commitment
    payload --- ten fields:
    \texttt{orderHash}, \texttt{processId}, \texttt{buyer},
    \texttt{seller}, \texttt{currency}, \texttt{payment},
    \texttt{cumulativeValue}, \texttt{agreementHash},
    \texttt{salt}, \texttt{deadline}. The \texttt{salt} and
    \texttt{deadline} fields are what enable the SDK's
    \texttt{reconstruct()} primitive to derive the EIP-712 digest
    deterministically from event logs alone.
  \item \texttt{OrderSeller} and \texttt{OrderCurrency} are
    companion events that index the seller and currency
    respectively---needed because the EVM caps indexed event
    arguments at three.
  \item \texttt{OrderResolved} fires once per order at resolution.
  \item \texttt{ProcessResolved} fires once per process at
    resolution, with the order count.
\end{itemize}

A reader who replays the four event types into a state machine
exactly reconstructs \texttt{processes}, \texttt{orderStatus}, and
\texttt{orderProcessId}---this is the basis of the SDK's
\texttt{reconstruct()} primitive used by agents and frontends.

% ════════════════════════════════════════════════════════════════════
\section{Design Non-Decisions}
\label{sec:non-decisions}

The following features are deliberately absent. Each absence
preserves a mechanism property
from~\citep{daliana2026mechanism}.

\begin{enumerate}
  \item \textbf{No owner, admin, or upgrade.} There is no
    \texttt{pause()}, \texttt{upgrade()}, or \texttt{transferOwnership()}.
    The contract is ownerless from deployment. Adding any of these would
    introduce an unbonded actor (the owner) into the resolution path,
    which the escape-hatch impossibility theorem rules out.
  \item \textbf{No protocol fee.} Settlement distributes the full
    bond amount to the parties. A fee at $k = 2$ would shift the
    buyer's net cooperation payoff above the defection payoff
    only by $\pay - \text{fee}$, narrowing the dominance margin and
    --- if the fee approaches $\pay$ --- collapsing the equilibrium.
  \item \textbf{No timeout.} An active order remains active
    indefinitely. Per the escape-hatch impossibility theorem, any
    timeout with recovery fraction $\alpha \geq \tfrac{1}{2}$ destroys
    weak dominance for the buyer.
  \item \textbf{No partial resolution.} The resolution function
    requires the full active order list and reverts otherwise. This
    enforces atomicity, on which the weakest-link coordination
    pressure depends.
  \item \textbf{No internal ledger.} Payouts are direct
    \texttt{safeTransfer} calls, not balance increments to be
    withdrawn later. This eliminates withdrawal-pattern reentrancy
    surface and removes a class of accounting drift bugs.
  \item \textbf{No restriction on buyer--seller equality.} The
    contract does not forbid $B = S$. If a single address signs both
    sides of a commitment, that address deposits both bonds and
    receives both payouts at resolution; the bond math is
    self-cancelling and no third party is exposed. We surface this
    explicitly because it is a reasonable design question to ask of
    every smart-contract escrow: in \texttt{FigaroCore}, the answer
    is that the equilibrium analysis treats $B \neq S$ as the
    standard case but the contract permits the degenerate case
    rather than introducing a guard whose effect would be cosmetic.
\end{enumerate}

\subsection{Custom Errors as Specification}

The kernel defines fifteen custom errors, each named for the
invariant it protects: \texttt{DeadlineExpired},
\texttt{InvalidBuyerSignature}, \texttt{InvalidSellerSignature},
\texttt{ZeroPayment}, \texttt{ProcessAlreadyExists},
\texttt{UnknownProcess}, \texttt{CumulativeValueMismatch(uint256
expected, uint256 actual)}, \texttt{NotProcessBuyer},
\texttt{CurrencyMismatch}, \texttt{OrderNotCommitted(bytes32
orderHash)}, \texttt{NoActiveOrders},
\texttt{IncompleteOrderList(uint256 required, uint256 provided)},
\texttt{DuplicateCommitment}, \texttt{FeeOnTransferDetected},
\texttt{InvalidRootCumulativeValue}.
The error names function as a partial executable specification:
each error is exercised by a Foundry revert-branch test and
referenced by name in the verification-results section below.

% ════════════════════════════════════════════════════════════════════
\section{Formal Verification Methodology}
\label{sec:verification-method}

We use four verification techniques, each chosen for what it covers
that the others do not.

\paragraph{TLA\textsuperscript{+} model checking (TLC).}
Captures the full state machine as a single transition system.
Invariants are propositional safety properties. TLC explores all
reachable states under bounded parameters by breadth-first search,
exhibiting either an invariant violation with a minimal trace or
exhaustive coverage of the bounded space. Strength: high-level
state-machine reasoning; abstracts cryptography and ERC-20 mechanics
to expose accounting errors. Weakness: bounded by parameters; gives
no guarantee for unbounded state.

\paragraph{Property-based fuzzing (Echidna).}
Runs randomized call sequences against the deployed bytecode under
HEVM, with EIP-712 signing performed via cheatcodes. Strength:
exercises real bytecode against concrete adversaries, including
combinations of operations the model abstracts. Weakness: random
sampling; finds bugs but does not prove their absence.

\paragraph{Symbolic execution (Halmos, z3-backed).}
Encodes call traces as SMT formulas and asks z3 whether any concrete
input satisfies the negation of an invariant. Strength: when it
returns \emph{verified}, the property holds for \emph{all} inputs in
the modeled trace; complements Echidna's bounded random sampling
with symbolic coverage of the commit/resolve state machine.
Weakness: symbolic complexity blows up with control-flow depth;
tractable on bounded path lengths.

\paragraph{Specification checking (Certora).}
SMT-based proving of CVL rules against the bytecode, run on
Certora's cloud service. Strength: rules can quantify over methods
(\texttt{method f}), allowing universal statements like ``no method
modifies $X$''. Weakness: paid service; rule encoding is non-trivial.

\paragraph{Foundry as the regression spine.}
A Foundry test suite (14 files, 225 tests as of snapshot date)
serves as the continuous-integration spine. Foundry tests are not
counted toward formal verification but provide the harness the other
three techniques compile against.

% ════════════════════════════════════════════════════════════════════
\section{Verification Results}
\label{sec:verification-results}

All numbers in this section are reported against the snapshot at
2026-04-23, repository commit
\texttt{970ee3a914b2}\footnote{Pin to this commit when reproducing
to guarantee identical bytecode; later commits may change call
counts as the test surface evolves.}. Re-runs produce different
fuzz-call counts; the configuration parameters (invariants, rule
counts, model bounds) are the stable artifact.

\subsection{TLA\textsuperscript{+} (\texttt{formal/FigaroCore.tla})}

The contract is modeled in TLA\textsuperscript{+} with eight state
variables (\texttt{processes}, \texttt{orderStatus},
\texttt{orderRecords}, \texttt{processOrders},
\texttt{contractBalance}, \texttt{wallets}, \texttt{nextProcId},
\texttt{nextOrdId}) and three actions: \texttt{CommitRoot},
\texttt{CommitSub}, \texttt{ResolveProcess}. The model abstracts
EIP-712 signatures (assumed correct given valid commitments), ERC-20
mechanics (modeled as integer balances), and timing (deadlines are
orthogonal to the bonding equilibrium and are exercised in Foundry
instead).

\paragraph{Model bounds (\texttt{formal/MC.cfg}).}
2 buyers ($b_1, b_2$), 2 sellers ($s_1, s_2$), initial balance 30,
maximum payment 3, 2 concurrent processes, 2 sub-orders per process.
Two buyers (rather than one) enable cross-buyer invariant
verification---a single-buyer model would not exercise
\texttt{NotProcessBuyer} branches reachable by an attacker buyer.

\paragraph{Invariants verified.}
Seven safety invariants verified exhaustively:
\begin{enumerate}[label=\textbf{I\arabic*}]
  \item \textbf{TypeOK}: type well-formedness of all state variables.
  \item \textbf{TokenConservation}: sum of wallet balances plus
    contract balance equals total supply.
  \item \textbf{ContractSolvency}: \texttt{contractBalance} $\geq 0$.
  \item \textbf{WalletNonNegative}: all participant balances $\geq 0$.
  \item \textbf{CumulativeIntegrity}: per-process,
    $\cumval$ equals the sum of all order payments.
  \item \textbf{ActiveCountCorrect}: stored active-order count
    matches the count of committed orders.
  \item \textbf{ResolutionAlwaysPossible}: for every active process,
    the contract holds sufficient funds to resolve all orders.
\end{enumerate}

\paragraph{Result.}
6{,}087{,}113 distinct states explored in 4 minutes 8 seconds; zero
invariant violations. Reproduced via \texttt{./test-tla.sh}.

\subsection{Echidna (\texttt{src/echidna/EchidnaFuzzerV5.sol})}

Seven property invariants, each prefixed \texttt{echidna\_}, checked
after every call sequence:
\begin{enumerate}[label=\textbf{E\arabic*}]
  \item \texttt{solvency}: token balance $\geq$ sum of outstanding
    bonds.
  \item \texttt{active\_count\_consistent}: stored count equals
    actual.
  \item \texttt{cumulative\_accounting}: accumulator equals sum of
    payments.
  \item \texttt{state\_monotonicity}: orders never transition
    backwards (0 $\to$ 1 $\to$ 2, never reverse).
  \item \texttt{token\_conservation}: total supply constant under
    all commit/resolve sequences.
  \item \texttt{buyer\_dominance}: non-buyer \texttt{resolveProcess}
    always reverts.
  \item \texttt{atomic\_resolution}: incomplete order list always
    reverts.
\end{enumerate}

\paragraph{Configuration (\texttt{echidna-v5.yaml}).}
\texttt{testLimit: 50000} (campaign budget), \texttt{seqLen: 30},
\texttt{shrinkLimit: 5000}, three deployer/sender accounts, property
mode, 300-second wall-clock timeout. The Echidna campaign budget is
the reproducibility anchor; raw call counts vary between runs.

\paragraph{Result.}
All seven properties held across the campaign; the most recent run
exercised approximately 43{,}000 individual calls before the budget
was exhausted, with no counterexample produced. Reproduced via
\texttt{./test-echidna.sh}.

\subsection{Halmos (\texttt{test/HalmosFigaroCore.t.sol})}

Seven kernel properties, each implemented as a \texttt{check\_}
function and discharged by z3 under symbolic-input quantification:
\begin{enumerate}[label=\textbf{H\arabic*}]
  \item \texttt{tokenConservation\_afterCommit}.
  \item \texttt{contractSolvency\_afterCommit}.
  \item \texttt{correctBondAmounts}: $\bbond = 2\pay$ and
    $\sbond = 2\cumval$ for all symbolic $\pay$.
  \item \texttt{resolutionPayouts}: $\pi_s = 2\cumval + \pay$ and
    $\pi_b = \pay$ for all symbolic $\pay$.
  \item \texttt{orderStatusTransition}: status moves only
    $0 \to 1 \to 2$.
  \item \texttt{buyerDominance\_revert}: any non-buyer caller of
    \texttt{resolveProcess} reverts.
  \item \texttt{cumulativeValueMonotonic}: across two symbolic
    sub-order commitments, $\cumval$ never decreases.
\end{enumerate}

\paragraph{Result.}
All seven discharged by z3. Reproduced via the Foundry profile
\texttt{halmos} as documented in the repo.

\subsection{Certora (\texttt{certora/})}

Five specification files. We list the rule counts per surface and
the headline rule for each.

\paragraph{Kernel (\texttt{FigaroCore.spec}, 8 rules).}
\texttt{orderStatusNeverDecreases},
\texttt{orderStatusTransitionsAreValid},
\texttt{commitIncreasesActiveCount},
\texttt{onlyBuyerCanResolve},
\texttt{noDoubleCommit},
\texttt{cumulativeValueMonotonic},
\texttt{rootBuyerImmutable},
\texttt{currencyImmutable}.

\paragraph{Attestation surface (\texttt{AttestationCoordinator.spec}, 9 rules).}
Six role-gate rules:
\texttt{nonBuyerCannotAttestAsBuyer},
\texttt{unknownProcessRevertsAsBuyer},
\texttt{successfulBuyerAttestationImpliesBuyer},
\texttt{buyerAttestationEnforcesProcessBoundary},
\texttt{attestationCannotChangeOrderStatus},
\texttt{attestationCannotChangeProcessState}.
Three validator-gate rules (verified on Certora cloud
2026-04-23):
\texttt{noValidatorBlocksBuyerAttestation},
\texttt{setValidatorIsFirstWriteWins},
\texttt{setValidatorPreservesOtherBindings}.

\paragraph{Token-operations conservation (\texttt{TokenOpsVerification.spec}, 7 declared rules).}
\texttt{commitBuyerExactDelta},
\texttt{commitSellerExactDelta},
\texttt{commitContractExactDelta},
\texttt{commitAllowanceDrainSafety},
\texttt{commitTokenConservation},
\texttt{resolveSingleOrderPayout},
\texttt{resolveSingleOrderConservation}.
A regression guard for ``no untracked token moves'' is documented in
the spec but is not itself a CVL rule. Certora's prover splits
\texttt{resolveSingleOrderConservation} into two sub-cases, so the
service report shows 8 sub-rules green for this spec.
The token-ops surface is gated by a pre-run lint
(\texttt{lint-token-ops.sh}, invoked from \texttt{test-certora.sh})
that asserts the spec's inventory file enumerates every ERC-20
\texttt{transfer}/\texttt{transferFrom}/\texttt{safeTransfer*} call
site in \texttt{src/}. New transfer call sites without a matching
inventory entry fail before any cloud dispatch.

\paragraph{Other specs.}
\texttt{StagedMerkleAirdrop.spec} (3 rules) covers the airdrop
contract, and \texttt{FigToken.spec} (6 rules) covers the FIG token
registry. Both are out of scope for the kernel claim of this paper
but ship in the same repository and are verified together.

\paragraph{Total Certora kernel-relevant rules.}
$8 + 9 + 7 = 24$ declared rules across three specs (Certora reports
25 sub-rules green when sub-case splits are counted). All
discharged on Certora cloud as of the snapshot date.

\subsection{Foundry Regression Spine}

14 test suites, 225 tests (including
\texttt{FigaroCoreRevertBranchTest} which targets every custom error
listed in \Cref{sec:non-decisions}; \texttt{ParityVectors} which
asserts EIP-712-digest parity between the SDK's TypeScript hash
function and the on-chain implementation; and
\texttt{FigaroCoreEventEmissionTest} which asserts every event field
encodes as the SDK expects).

\subsection{Scope of the Verification Claim}
\label{sec:verification-scope}

What this paper's verification establishes, in plain language for
non-specialist readers: the Solidity source at commit
\texttt{970ee3a914b2} has been audited by four different
methodological approaches, and each approach has found that the
source satisfies seven named properties within the bounds each
approach explores. What this does \emph{not} establish:

\begin{itemize}
  \item \textbf{That the specification is the right specification.}
    Verification checks code against a spec; it does not check
    whether the spec captures what the author meant or what users
    expect.
  \item \textbf{That bytecode at a deployed address was compiled
    from this commit.} The kernel's production deployment is a
    separate act, performed by a separate party, using a compiler
    under settings and dependency versions that this paper does
    not pin. Reproducing the verified bytecode against a deployed
    address requires standard contract-source-verification
    tooling (Etherscan-style verification, sourcify, or
    reproducible-build pipelines) which is out-of-scope here.
    Relying parties should confirm the source-to-bytecode-to-address
    chain independently.
  \item \textbf{That the denomination token behaves.} The kernel
    takes an arbitrary \texttt{IERC20}. If the chosen token's
    issuer exercises blocklist, freeze, upgrade, or pause
    authority over a party's bonded balance, the ``no escape
    hatch'' property of the kernel is vacuously violated --- the
    bond is still locked from the kernel's perspective, but it
    is unreachable from the world's. Parties should treat the
    choice of denomination token as a selection of whose escape
    hatches to accept, not as an avoidance of escape hatches in
    aggregate.
  \item \textbf{That chain liveness and key custody persist.} If
    the host chain halts, censors, or reorgs past finality, and
    if either party loses or has compromised their signing key
    during the process, the kernel has no recovery mechanism.
    These are operational-security dependencies that sit above
    the verified surface.
\end{itemize}

The honest summary is that verification reduces one class of risk
(specification--implementation gap within the verified bounds) and
displaces several others (specification correctness, deployment
chain of custody, token-issuer authority, host-chain liveness, key
custody) onto layers the paper does not itself verify. A
practitioner or legal reader consuming these verification claims
should treat them as input to due diligence, not as a substitute
for it.

% ════════════════════════════════════════════════════════════════════
\section{Security Analysis}
\label{sec:security}

\subsection{Threat Model}

We assume rational adversaries who maximize economic payoff. We
consider four attack vectors.

\paragraph{Griefing (buyer refuses to resolve).}
The attacker's cost is $2\pay$ (locked bond) plus opportunity cost on
that capital plus permanent on-chain reputation damage. The attacker's
benefit is zero. For rational agents, the attack is strictly
dominated by cooperation. For irrational agents (spite exceeds
economic loss), Layer~3 (below) provides deterrence: the immutable
on-chain record serves as evidence in off-chain legal proceedings.

\paragraph{Sybil (fake orders to manipulate state).}
Each order requires real bond deposits totaling $2(\cumval_i +
\pay_i)$ per order. At scale, Sybil attacks are capital-intensive
with no path to extracting the invested capital---there is no
secondary market for locked bonds, no yield, no governance influence.

\paragraph{Front-running.}
Order identifiers are content-addressed from the dual-signed
commitment. An attacker who observes a pending \texttt{commit}
transaction cannot create a conflicting order for the same process:
producing a different signature changes the digest, which changes
the order hash, so the front-runner's order would be a distinct
order, not a conflicting one. No extractable MEV exists.

\paragraph{Cumulative value manipulation.}
Per the self-enforcing-DAG-honesty result
in~\citep{daliana2026mechanism}, overstating cumulative value is
strictly dominated (greater capital at risk, same payoff).
Understating is prevented by the contract's monotonic accumulator
check (\texttt{CumulativeValueMismatch} revert).

\subsection{Liveness}

\begin{theorem}[Liveness Under Rationality]
\label{thm:liveness}
Let $B$ be a buyer with time-discounted utility
$U_B = \sum_{t=0}^{\infty} \delta^t \cdot u_t$ where
$\delta \in (0, 1)$ is the discount factor and $u_t$ is the
per-period payoff. If all sellers have cooperated (the buyer has
received the agreed goods), then the buyer resolves in finite time.
\end{theorem}

\begin{proof}
The buyer's payoff under \emph{resolve at $t^*$} is $\delta^{t^*}
\sum_i \pay_i$; under \emph{never resolve} is $0$. Since
$\delta \in (0,1)$ and $\pay_i > 0$, every finite $t^*$ strictly
dominates non-resolution; the optimum is $t^* = 0$.
\end{proof}

\begin{remark}[Limitations]
\label{rem:liveness-limitations}
\Cref{thm:liveness} assumes (i)~the buyer received satisfactory
performance, (ii)~$\delta < 1$, and (iii)~no spite. Condition (iii)
is the hardest to guarantee. When it fails, the mechanism cannot
guarantee liveness by economic incentives alone---the residual case
is addressed by Layer~3 (legal deterrence via immutable on-chain
evidence).
\end{remark}

\subsection{The Three-Layer Enforcement Architecture}

\begin{table}[ht]
\centering
\caption{Defense-in-depth: three enforcement layers.}
\label{tab:enforcement}
\begin{tabular}{llll}
\toprule
\textbf{Layer} & \textbf{Mechanism} & \textbf{Coverage} & \textbf{Failure Mode} \\
\midrule
1. Bonding & Asymmetric collateral & Rational actors & Irrational actors \\
2. Coordination & Atomic resolution & Multi-seller failures & Uncoordinated sellers \\
3. Evidence & Immutable on-chain log & Adversarial actors & Jurisdictional limits \\
\bottomrule
\end{tabular}
\end{table}

\paragraph{Layer 1 (economic).}
The bonding equilibrium of~\citep{daliana2026mechanism} handles
rational participants; cooperation strictly dominates defection at
every node.

\paragraph{Layer 2 (coordinational).}
Atomic resolution induces the weakest-link subgame; honest sellers
have direct economic interest of magnitude $\pay_i + 2\cumval_i$ in
ensuring co-sellers cooperate. No explicit communication is required.

\paragraph{Layer 3 (evidentiary).}
Every commitment, every bond deposit, every resolution (or its
absence) is recorded with block timestamps. The kernel does not
attempt on-chain dispute resolution; it provides the \emph{evidence}
that off-chain dispute systems already accept.

\subsection{Scope and Limitations}

\paragraph{Liveness under non-rational failure.}
Catastrophic non-rational failure (lost keys, hardware failure,
participant death) imposes losses on co-sellers via atomic
resolution regardless of incentive alignment. The
coordination-pressure result assumes payoffs depend on co-sellers'
\emph{choices}; failure is not a choice. Operational reliability
under such failure depends on participant continuity practices (key
management, multisig, delegated signing) that sit above the kernel
and are out of scope.

\paragraph{Capital opportunity cost.}
At prevailing risk-free rate $r$ over process duration $t$, the
buyer and seller forgo $r\pay t$ and $2r\cumval t$ respectively.
Dominance is preserved at any $rt > 0$, but the operational margin
narrows; the $2\times$ multiplier is sufficient as a matter of
game-theoretic derivation, but the practical applicability favors
processes short relative to the prevailing risk-free horizon.

% ════════════════════════════════════════════════════════════════════
\section{Mechanism Composition: The Coordinator Pattern}
\label{sec:composition}

The kernel's bonding equilibrium is a fixed point. The natural next
question---under what conditions does an external mechanism composed
with the kernel preserve the equilibrium?---is answered, in this
implementation, by the \emph{coordinator pattern}.

\begin{proposition}[Coordinator Pattern --- Sufficient Conditions]
\label{prop:coordinator}
An external mechanism $M$ composed with \texttt{FigaroCore} preserves
the bonding equilibrium of~\citep{daliana2026mechanism} if it
satisfies:
\begin{enumerate}[label=(\roman*)]
  \item \textbf{Read-only kernel access}: $M$ interacts with kernel
    state via \texttt{staticcall} or event consumption only.
  \item \textbf{Role non-impersonation}: $M$ never calls
    \texttt{commit} or \texttt{resolveProcess} on behalf of a buyer
    or seller; signatures are originated by the parties themselves.
  \item \textbf{Bond non-modification}: $M$ holds no kernel bonds
    and does not interpose between the parties and the bond
    transfers.
  \item \textbf{No bypass}: $M$ does not provide an alternative
    resolution path or escape from \texttt{Active}.
\end{enumerate}
\end{proposition}

\Cref{prop:coordinator} is sufficient, not necessary; mechanisms
that violate one of (i)--(iv) may still preserve the equilibrium
under additional argument, but those arguments are per-mechanism.
The four conditions are currently checked by inspection per
extension contract rather than by a parametric Certora rule that
quantifies over arbitrary $M$ --- see \Cref{sec:conclusion} for the
verification status of this proposition. The envelope is the
discipline that makes composition \emph{by construction} safe.
We exhibit three instances.

\paragraph{Dutch auction (\texttt{src/DutchAuction.sol}).}
A descending-price coordination primitive that performs role
allocation: who may become the seller for a given root commitment.
The auction holds no tokens, custodies no bonds, and never calls
\texttt{commit}. The auction's winner subsequently signs an
EIP-712 commitment off-chain and the parties execute \texttt{commit}
themselves. Conditions (i)--(iv) hold by construction.

\paragraph{Attestation coordinator (\texttt{src/AttestationCoordinator.sol}).}
A unified zero-storage attestation surface, validator-gated per
\texttt{schemaId}. Three modes: \texttt{attestAsSeller},
\texttt{attestAsBuyer}, \texttt{attestViaResolver}. Each call
verifies role membership against the kernel's
\texttt{ProcessState.rootBuyer} (or against an
\texttt{IRoleResolver} for the resolver path), runs
\texttt{schemaValidator[schemaId].validate(...)} on the content, and
emits an \texttt{Attestation} event whose \texttt{contentRef} is
$\mathrm{keccak256}(\text{content})$. The coordinator never modifies
kernel state---this is the headline Certora rule
\texttt{attestationCannotChangeOrderStatus} and its sibling
\texttt{attestationCannotChangeProcessState}, both verified
exhaustively. Conditions (i)--(iv) hold; the coordinator is the
cleanest example of an event-only extension.

\paragraph{Schema and operator registries (\texttt{src/SchemaRegistry.sol}, \texttt{src/OperatorRegistry.sol}).}
Permissionless, append-only event-first registries. The schema
registry anchors content schemas as \texttt{keccak256(name)} with a
URI hash pointing at the off-chain JSON spec. The operator registry
is on-chain operator self-registration with a reclaimable ETH
deposit. Both compose with the kernel by being visible to readers
and not being callable by \texttt{commit} or
\texttt{resolveProcess}. Conditions (i)--(iv) hold; both registries
are pure coordination surfaces with no kernel interaction at all.

\paragraph{What this gives.}
\Cref{prop:coordinator} is the discipline used to evaluate every
proposed extension. A mechanism that satisfies all four conditions
can be composed without re-doing the equilibrium proof; a mechanism
that violates one of them requires explicit mechanism-level argument
and re-verification. In the reference implementation, all five
extension contracts shipped to date satisfy
\Cref{prop:coordinator} by construction.

\paragraph{A note on liability allocation when conditions fail.}
The verification status of \Cref{prop:coordinator} is important
for legal readers: the four conditions are presently checked by
inspection per extension contract rather than by a parametric
Certora rule that quantifies over arbitrary $M$ (see
\Cref{sec:conclusion} for the verification gap this leaves open).
If an extension contract silently violates condition (i), (ii),
(iii), or (iv) and a downstream user is harmed when the kernel
equilibrium breaks for that composition, the kernel verification
in this paper does not extend to the resulting harm. Liability in
such a case is allocated by the legal regime governing the
extension deployer's conduct, the audit artifacts (or their
absence) accompanying the extension, and the disclosures made to
users --- all of which sit outside the kernel and are treated in
the law-and-evidence companion~\citep{daliana2026legal}. The
kernel claims it is correct; it does not claim that every
extension built against it is correct, nor does it claim that
the legal chain from kernel correctness to user remedy is
closed.

% ════════════════════════════════════════════════════════════════════
\section{Conclusion and Future Work}
\label{sec:conclusion}

We have presented a verified Solidity implementation of the
asymmetric-bonding mechanism of~\citep{daliana2026mechanism}: an
ownerless, fee-less kernel of two functions and three mappings, with
a four-layer verification stack (TLA\textsuperscript{+}, Echidna,
Halmos, Certora) that exercises the same invariants from four
methodological directions. The coordinator pattern of
\Cref{sec:composition} provides a discipline for extending the
kernel without weakening its equilibrium.

\paragraph{Agent-native operation.}
The kernel is designed for autonomous agent participation as a
first-class concern. EIP-712 typed commitments allow agents to
build, verify, and sign structured data using standard libraries.
Event-sourced state reconstruction means agents can derive the
complete process graph from on-chain logs without subgraph
dependencies. The companion SDK exposes a typed action proposer
(\texttt{proposeActions}) and a dual-mode execution queue
(\texttt{ActionQueue}) that supports both human-in-the-loop approval
and fully autonomous submission. Nothing in the kernel
distinguishes a human signer from an autonomous agent: the bonding
equilibrium holds for any rational actor capable of controlling an
address and bearing capital risk.

\paragraph{Execution-environment portability.}
A Cairo/Starknet port of the kernel is in progress. The mechanism is
execution-environment-agnostic; what varies between EVM and Cairo is
the cryptographic primitive (ECDSA vs.\ STARK signatures), the
storage cost model, and the event surface. The verification stack
must be re-instantiated per environment: TLA\textsuperscript{+} and
Certora rules are mostly portable (state-machine reasoning); Halmos
is EVM-specific; Echidna ditto. The conformance specification
between EVM and Cairo is the SP1 prover's Rust kernel, which mirrors
the TS validator and Solidity hashing byte-for-byte.

\paragraph{Verification gaps the present snapshot does not close.}
The Echidna corpus is accumulated but not minimized; corpus
maintenance is a continuous task, not a release event.
\Cref{prop:coordinator} is currently asserted at the implementation
level (the four conditions are checked by inspection per extension
contract); a formal Certora rule that checks
``$M$ does not call \texttt{commit}'' for arbitrary $M$ has not yet
been authored. A future revision of this snapshot will include it,
along with the reverse direction (no kernel storage write reachable
from any extension entry point).

\paragraph{The deeper claim: defense in depth, not a probability bound.}
Beyond the headline numbers, the value of the verification stack
is that four methodologically different proofs converge on the
same seven properties (token conservation, solvency, bond
correctness, status monotonicity, buyer dominance, cumulative
monotonicity, atomic resolution). We present this as
\emph{defense in depth}: a bug that survives all four toolchains
must be consistent with the bounded-exhaustive exploration of
TLA\textsuperscript{+}, the property-level random sampling of
Echidna, the symbolic-input discharge of Halmos, and the
method-quantified SMT rules of Certora. We do not attempt a
formal probability bound on residual error: Halmos and Certora
both lean on the z3 SMT solver, which breaks the independence
assumption any probabilistic bound would require, and the prior
distribution over kernel-invariant violations in this toolchain
ecosystem is not characterized in the literature. The honest
statement is that four methodologically distinct proofs agreeing
provides substantially stronger assurance than any one alone, not
that the residual probability is bounded above by any particular
value. The stack reduces risk; it does not retire it.

\section*{Acknowledgments}
The verification surface described here owes its current shape to
many hands. I want in particular to acknowledge Vitalik Buterin for
the Safe Remote Purchase primitive that seeded the work, Vlad Zamfir
for the mechanism-design lens that made the generalization legible,
Fabrizio Romano Genovese for his contributions to \figaro{}, the
Coalition of Automated Legal Applications (COALA), and the
maintainers of TLA\textsuperscript{+}, Echidna, Halmos, and Certora.


% ════════════════════════════════════════════════════════════════════
% References
% ════════════════════════════════════════════════════════════════════

\begin{thebibliography}{99}

\bibitem[Daliana(2026a)]{daliana2026mechanism}
Daliana, A.
\newblock Asymmetric Bonding: A Self-Enforcing Settlement Primitive
  for $N$-Party Coordination.
\newblock Companion mechanism-design paper, 2026.

\bibitem[Daliana(2026b)]{daliana2026legal}
Daliana, A.
\newblock On-Chain Evidence, Off-Chain Adjudication: A Layered
  Approach to Smart-Contract Dispute Resolution.
\newblock Companion law-and-evidence paper, 2026.

\bibitem[Daliana(2026c)]{daliana2026fig}
Daliana, A.
\newblock \textsc{Fig}: A Coordination Token Decoupled from
  Protocol Activity.
\newblock Companion cryptoeconomics paper, 2026.

\bibitem[EIP-712(2017)]{eip712}
Bloemen, R., Logvinov, L., and Evans, J.
\newblock {EIP-712}: Ethereum typed structured data hashing and signing.
\newblock Ethereum Improvement Proposal, 2017.
\newblock \url{https://eips.ethereum.org/EIPS/eip-712}

\end{thebibliography}

\end{document}
